%=============================================================================!
% 14.S  WHAT WE NOW HAVE                                                      !
%=============================================================================!
\unnumbered{What we now have}

The pressure of a gas is the momentum its atoms deliver to the walls per
unit time and area, $N\kB T/V$ for atoms that do not interact.
Equipartition applied to the positions adds the virial's share, $PV =
N\kB T + \frac13\ens{\mathcal{W}}$, and the free energy's response to a
change of volume gives the same formula in a periodic box, with the virial
summed over pairs from their separations to the nearest copy; the sum
$\sum_i\vb{r}_i\cdot\vb{F}_i$ from wrapped positions changes with the
origin.

Its response to a strain gives the pressure tensor,
$(\sum_im_iv_{i\alpha}v_{i\beta} + \mathcal{W}_{\alpha\beta})/V$, isotropic for
a liquid, uneven for a strained solid; ASE and many libraries report the
stress, its negative.

A barostat must sample $\mathcal{P}(V) \propto
\mathrm{e}^{-\beta P_0V}Z(N, V, T)$, which for an ideal gas is
$V^N\mathrm{e}^{-\beta P_0V}$ and for a liquid has the variance $\kB
T\ens{V}\kappa_T$.

Berendsen's barostat holds the mean pressure but squeezes the
volume's fluctuations, which then give less than half the liquid's
compressibility; stochastic cell
rescaling adds to $\ln V$ the noise the fluctuation-dissipation relation
demands and samples the volume correctly; the MTK piston gives the volume
a mass, swings with the period $2\pi\tau_{\mathrm P}\sqrt{(N + 1)\kB T
\kappa_T/3V}$, and with chains samples correctly too.

A solid that wants to
change shape needs semi-isotropic or anisotropic coupling: isotropic
coupling of a layered solid with guests between its sheets leaves it
pulled along the sheets and squeezed across them while its mean pressure
reads the target.

A healthy run shows a widely scattered instantaneous
pressure whose running mean settles at $P_0$, a volume that spreads by
$\sqrt{\kB T\ens{V}\kappa_T}$, and an effective energy that holds still.

Chapter~15 asks how long a run must be before such averages, and the
error bars on them, can be trusted.
